% Akshat Sachin Valse — curriculum vitae
% Compile:  latexmk -pdf cv.tex     Clean: latexmk -c
% Output is copied to ../public/<PDF_FILE> (dated filename, see src/pages/cv.astro);
% /cv.pdf redirects to the current file via public/_redirects.

\documentclass[11pt]{article}

\usepackage[margin=0.85in,top=0.7in,bottom=0.65in]{geometry}
\usepackage{amsmath,amssymb}
% Latin Modern rather than bitmap CM: same design, scalable outlines, which
% microtype's font expansion requires and which keeps the PDF text selectable.
\usepackage{lmodern}
\usepackage[T1]{fontenc}
\usepackage{microtype}
\usepackage{enumitem}
\usepackage{titlesec}
\usepackage[hidelinks]{hyperref}
\usepackage{fancyhdr}

\hypersetup{
  pdftitle={Akshat Sachin Valse -- Curriculum Vitae},
  pdfauthor={Akshat Sachin Valse},
  pdfsubject={Curriculum vitae},
  pdfkeywords={stochastic analysis, Langevin sampling, large deviations, rare-event simulation, neural PDE solvers},
}

\pagestyle{fancy}
\fancyhf{}
\renewcommand{\headrulewidth}{0pt}
\fancyfoot[R]{\footnotesize Valse \textperiodcentered\ \thepage}
\fancyfoot[L]{\footnotesize Last updated \today}

\titleformat{\section}{\normalsize\scshape}{}{0pt}{}[\vspace{-0.6em}\rule{\textwidth}{0.4pt}]
\titlespacing{\section}{0pt}{0.85em}{0.4em}

\setlist[itemize]{leftmargin=1.2em,itemsep=0.15em,topsep=0.25em,parsep=0pt}

% Entry with a right-aligned date. Two minipages rather than \hfill so a long
% title wraps in its own column instead of running into the date.
\newcommand{\entry}[2]{%
  \noindent\begin{minipage}[t]{0.78\textwidth}\raggedright\textbf{#1}\end{minipage}%
  \hfill%
  \begin{minipage}[t]{0.22\textwidth}\raggedleft #2\end{minipage}\par\vspace{0.1em}}
\newcommand{\subentry}[1]{\noindent\textit{#1}\par\vspace{0.15em}}
\newcommand{\course}[2]{#1 (\textsc{#2})}

\begin{document}

\begin{center}
  {\LARGE Akshat Sachin Valse}\\[0.5em]
  {\small
    Department of Mathematics, Iowa State University, Ames, Iowa\\
    \href{mailto:avalse@iastate.edu}{avalse@iastate.edu}
    \ \textperiodcentered\ \href{https://akshatvalse.com}{akshatvalse.com}
    \ \textperiodcentered\ \href{https://github.com/AkshatValse}{github.com/AkshatValse}
  }\\[0.6em]
  {\small\itshape Stochastic analysis for machine learning: Langevin sampling and its
  discretization,\\ large deviations and rare-event simulation, and the stability of
  learned PDE solvers.}
\end{center}

% ---------------------------------------------------------------------------
\section{Education}

\entry{Iowa State University, Honors Program}{Expected May 2028}
\subentry{B.S.\ Statistics \textperiodcentered\ B.S.\ Mathematics \textperiodcentered\
B.A.\ Computer Science \textperiodcentered\ B.S.\ Data Science}
\noindent Cumulative GPA 3.85 / 4.00.

\vspace{0.5em}
\noindent\textbf{In progress, Fall 2026} (all five at the PhD level, 6000-series).
\course{Foundations of Probability Theory, measure-theoretic}{stat 6541};
\course{High-Dimensional Probability and Linear Algebra for Machine Learning}{math 6230};
\course{Advanced Topics in Machine Learning}{com s 6730};
\course{Advanced Topics in Artificial Intelligence}{com s 6720};
\course{Advanced Topics in High-Performance Computing}{com s 6250}.

\vspace{0.35em}
\noindent\textbf{Planned, Spring 2027.} Advanced Probability Theory; Advanced Stochastic
Processes.

\vspace{0.35em}
\noindent\textbf{Graduate and advanced coursework, completed.}\par
\vspace{0.15em}
\noindent\textit{Probability and statistics:}
\course{Theory of Probability and Statistics}{stat 5430};
\course{Stochastic Process Models}{stat 5540};
\course{Empirical Methods for the Computational Sciences}{stat 5830};
\course{Multivariate Analysis}{stat 5750};
\course{Categorical Data Analysis}{stat 5770};
\course{Statistical Computing}{stat 5860}.\par
\vspace{0.15em}
\noindent\textit{Mathematics:}
\course{Numerical Methods for Differential Equations}{math 5810};
\course{Analysis I}{math 4140};
\course{Introduction to Mathematical Logic}{math 4430x};
\course{Partial Differential Equations}{math 3850};
\course{Theory of Linear Algebra, honors}{math 3170}.\par
\vspace{0.15em}
\noindent\textit{Computing and machine learning:}
\course{Concepts and Applications of Machine Learning, honors}{ds 3030};
\course{Design and Analysis of Algorithms}{com s 3110};
\course{Software Development Practices}{com s 3090}.

% ---------------------------------------------------------------------------
\section{Research Experience}

\entry{Langevin Sampling and Diffusion Estimation for Molecular Systems}{May 2026 -- present}
\subentry{Research Collaborator, Department of Mathematics, Iowa State University.
Advisor: Prof.\ David Herzog. Funded by NSF award DMS-2246491.}
\begin{itemize}
  \item Built an end-to-end Python simulation pipeline implementing MALA and projected
    Euler--Maruyama samplers for overdamped Langevin dynamics, staged from a
    one-dimensional periodic potential to a solvated ion in a three-dimensional
    Lennard-Jones bath; estimated self-diffusion and mobility coefficients from
    long-run trajectory data.
  \item Derived the closed-form Lifson--Jackson diffusivity
    $D = 1/I_0(1)^2 \approx 0.62386$ as an exact benchmark and validated the estimators
    against it before scaling to three dimensions, where the pipeline reproduces
    published solvated-ion results.
  \item Author weekly technical reports, including a companion note supporting the
    group's CLT-based error analysis.
\end{itemize}

\vspace{0.35em}
\entry{Large Deviation Theory for Rare-Event Analysis in Queueing Systems}{Nov.\ 2024 -- present}
\subentry{Undergraduate Research Fellow, Department of Mathematics, Iowa State
University. Advisor: Prof.\ Ruoyu Wu. Dean's High Impact Undergraduate Research Award.}
\begin{itemize}
  \item Developed the first importance-sampling methodology for rare-event estimation
    in the mean-field join-the-shortest-queue model, deriving state-dependent
    exponential tilting from the Budhiraja--Friedlander--Wu (2021) large-deviation rate
    function; stated a logarithmic-efficiency conjecture with a proof sketch via the
    Dupuis--Wang subsolution framework.
  \item Built a reproducible, performance-tuned Python simulation framework
    (Xoshiro256** with deterministic seeding) and validated it across 27
    configurations ($n \le 100$, $\rho \in \{0.7, 0.9, 0.95\}$): estimates agree with
    crude Monte Carlo within 95\% confidence intervals wherever crude Monte Carlo
    succeeds, and extend to probabilities as small as $10^{-31}$, where it fails
    entirely.
  \item Designed and benchmarked two heuristic tilting families (rate-reversal,
    level-weighted) with pilot-run selection; diagnosed over-tilting at extreme rarity
    through likelihood-ratio calibration statistics, motivating the ongoing computation
    of the optimal state-dependent tilt.
\end{itemize}

\vspace{0.35em}
\entry{Stability Analysis of Cell-Average Neural Networks for PDEs}{May 2026 -- present}
\subentry{Undergraduate Researcher, Department of Mathematics, Iowa State University.
Advisor: Prof.\ Jue Yan. Funded by the NSF D4 NRT. Joint with Mengyuan Yang.}
\begin{itemize}
  \item Adapting von Neumann stability analysis to neural-network PDE solvers: deriving
    the linearized amplification factor of the learned time-stepping scheme and
    characterizing how perturbations propagate through it, toward stability conditions
    on the trained weights analogous to those of classical finite-difference methods.
  \item Reimplementing the group's solver in Python to enable systematic stability
    experiments across schemes and step sizes.
\end{itemize}

% ---------------------------------------------------------------------------
\section{Manuscripts in Preparation}

\begin{enumerate}[leftmargin=1.4em,itemsep=0.4em,topsep=0.3em]
  \item \textit{Large-Deviation-Guided Importance Sampling for Rare Events in the
    Mean-Field Join-the-Shortest-Queue Model.} With R.\ Wu.

  \item \textit{Langevin Dynamics and Transport-Coefficient Estimation for Molecular
    Systems.} With D.\ Herzog.

  \item \textit{Stability Certificates for Cell-Average-Based Neural Network Solvers:
    von Neumann Analysis of Trained Flux Networks.} With M.\ Yang and J.\ Yan.
\end{enumerate}

% ---------------------------------------------------------------------------
\section{Awards and Honors}

\begin{itemize}
  \item Dean's High Impact Undergraduate Research Award, Iowa State University, 2026.
  \item Schillmoeller Family Scholarship in Statistics.
  \item Mead Fund Award in Statistics.
  \item James L.\ Cornette Mathematics and Interdisciplinary Studies Award.
  \item Dean's List, every semester.
\end{itemize}

% ---------------------------------------------------------------------------
\section{Leadership and Service}

\begin{itemize}
  \item Treasurer, Iowa State Data Science Club, 2025 -- present. Secured \$60{,}000+
    in funding supporting 130+ members.
  \item Research Ambassador, Iowa State University Honors Program.
  \item FAA Private Pilot Certificate, high-performance endorsement, 2025. Volunteer
    pilot and recruitment advisor, Pilots N Paws animal-rescue flights.
\end{itemize}

% ---------------------------------------------------------------------------
\section{Technical Skills}

\noindent\textit{Languages and tools.} PyTorch and CUDA for neural-network and
GPU-accelerated simulation; Python (NumPy, SciPy, pandas, scikit-learn, matplotlib);
R; \LaTeX; git.\par
\vspace{0.3em}
\noindent\textit{Methods.} Measure-theoretic probability and stochastic differential
equations; MCMC (MALA, Metropolis--Hastings) and importance sampling; large-deviation
asymptotics; numerical analysis of time-stepping schemes and von Neumann stability;
reproducible computation with recorded seeds and parameter sidecars, so every figure
regenerates from saved results.

\end{document}
